% §19.1 What This Chapter Is For
%
% Monograph register: motivates the second lens. The whole book has developed
% the framework set-theoretically (membership predicate, FPR/FNR, set algebra);
% this chapter re-reads the same objects as vectors and stochastic matrices, a
% parallel and complementary view, not a replacement. The operator view exposes
% the linear and spectral structure (composition = product, independence =
% tensor factorization) the set view leaves implicit; the set view keeps the
% combinatorial semantics the matrix view flattens. This section also fixes the
% local bra-ket / tensor notation used throughout the chapter (defined here, not
% in alex.sty, because the chapter is the only consumer). Cross-refs
% ax:element-indep, sec:kronecker, sec:composition-theorem, sec:n-ary-channels,
% sec:quantum-contextuality. DEFINED: sec:operator-what-for.
%
% Local notation macros for this chapter (bra-ket and a wider Kronecker spacing).
% \providecommand keeps the book safe if these names appear elsewhere later.
\providecommand{\ket}[1]{\lvert #1 \rangle}
\providecommand{\bra}[1]{\langle #1 \rvert}
\providecommand{\ketbra}[2]{\lvert #1 \rangle\!\langle #2 \rvert}
\providecommand{\Kron}{\mathbin{\otimes}}

\section{What This Chapter Is For}
\label{sec:operator-what-for}

Every preceding chapter developed the framework in one voice. A set is a
membership predicate; an approximate set is that predicate made noisy; the
errors are the rates $\fprate$ and $\fnrate$; the operations are union,
intersection, and complement, with rates that compose by the law of total
probability. That development is complete on its own terms, and nothing in this
chapter revises it.

This chapter reads the same objects in a second voice. Each element of the
universe becomes a basis vector, a set becomes a vector built from those, a
function becomes a matrix, and a noisy channel becomes a stochastic matrix
acting on vectors. The bra-ket and tensor notation of finite-dimensional linear
algebra carries the account. The two voices describe one framework, and the
point of stating both is that each makes visible what the other leaves
implicit.

\subsection*{What the operator view exposes}

The set-theoretic development states the composition theorem
(\Cref{sec:composition-theorem}) and the Kronecker factorization
(\Cref{sec:kronecker}) as two separate results, proved by two separate
arguments, one about chained gates and one about parallel queries. In the
operator view they are one fact read in two geometries: serial composition is a
matrix product and parallel querying is a tensor product, and both are products
of operators. The composition theorem's $1 - \prod_i (1 - \eta_i)$ form is then
not a combinatorial identity to be rediscovered but the multiplicativity of a
product, and the spectral structure that the product form hides becomes
available: eigenvalues multiply along a chain, and the error of a parity over
independent bits is a product of the channels' eigenvalues. The independence
axiom (\Cref{ax:element-indep}) acquires a one-line restatement, ``the noise
operator factors,'' and parametric parsimony becomes the statement that a
factored operator carries the sum of its factors' parameters, not their
product. None of this is new content; it is the same content with its linear and
spectral structure made explicit.

\subsection*{What the set view keeps}

The trade runs the other way as well. The matrix view flattens the
combinatorial semantics that the set view carries on its surface. A
false-positive rate written as a matrix entry $Q_{yx}$ no longer announces that
it is the probability of accepting a non-member; a union written as an
elementwise OR on occupation vectors no longer displays the three truth-regions
($A \setminus B$, $B \setminus A$, outside both) that
\Cref{sec:other-set-ops} weights to get the intersection rate right. The
operator view can compute these quantities, but it does not keep them legible
the way the predicate language does. Where the combinatorial meaning is the
point, the set view is the better instrument, and the matrix is the supporting
calculation.

\subsection*{One framework, two lenses}

So the relationship is parallel and complementary, not hierarchical. A reader
who wants the meaning of an operation reaches for the predicate; a reader who
wants its algebraic or spectral structure reaches for the operator. The
remaining sections build the operator lens in order: the computational basis
(\Cref{sec:operator-computational-basis}), maps as stochastic matrices
(\Cref{sec:operator-maps-stochastic}), the subtler case of sets as tensored
channels (\Cref{sec:operator-sets-tensored}), the algebra and spectra that the
view exposes (\Cref{sec:operator-composition-spectra}), and a precise account
of how far the formalism coincides with finite-dimensional quantum mechanics
and where, at \Cref{sec:quantum-contextuality}, the coincidence stops
(\Cref{sec:operator-classical-sector}). The bra-ket notation throughout is a
notational convenience for the classical stochastic algebra, adopted because it
makes the products and factorizations easy to write, and it is not a claim that
the framework is quantum mechanical.
